\subsection{Medical Retrieval-Augmented Question Answering}

Medical RAG systems~\citep{xiong2024benchmarking} have shown that retrieval quality dominates downstream QA performance, often mattering more than the choice of generator model. Iterative medical RAG~\citep{xiong2024imedrag} tries to expand evidence coverage by generating follow-up retrieval queries. Benchmarking studies systematically evaluate corpora, retrievers, and backbone models~\citep{xiong2024benchmarking}. But every system in this space shares a basic architecture: each passage gets scored independently against the query. Evidence selection stays pairwise. None of these pipelines has a mechanism for reranking evidence using the higher-order biomedical relations that connect passages to each other and to structured knowledge. A passage that shares a MeSH ancestor with the query but no surface lexical overlap remains invisible no matter how many retrieval rounds you run.

\subsection{Biomedical Semantic Retrieval}

MedCPT~\citep{jin2023medcpt} trains a dual-encoder retriever and a cross-encoder reranker on large-scale PubMed search logs. These are strong neural baselines that capture biomedical semantics far better than generic retrievers. The BioASQ benchmark~\citep{tsatsaronis2015bioasq,krithara2023bioasqqa} provides the standard evaluation framework. PubMedQA~\citep{jin2019pubmedqa} and BEIR NFCorpus~\citep{thakur2021beir} offer complementary testbeds. All of these semantic models operate on isolated query-passage pairs. They learn from the statistical patterns of PubMed click logs. They do not use freely available structured biomedical knowledge: MeSH hierarchies, biomedical entity graphs, curated relation databases like PrimeKG. Our question is whether adding these knowledge sources as ranking features improves evidence reranking beyond what semantic models alone achieve.

\subsection{Graph and Hypergraph Approaches to RAG}

A growing number of systems apply graph structure to retrieved information. GraphRAG~\citep{edge2024graphrag} organizes passages into entity-relationship graphs and performs local-to-global search for summarization. Medical Graph RAG~\citep{wu2024medgraphrag} grounds medical responses by linking documents to credible sources and controlled vocabularies. HyperGraphRAG~\citep{luo2025hypergraphrag} uses hyperedges for passage clustering and argues that n-ary relations capture complex information structures more naturally than pairwise edges. Cross-granularity HGRAG~\citep{wang2026hgrag} models entities and passages at different retrieval granularities. Medical HyperRAG~\citep{xu2026medicalhyperrag} builds a ClinBridge hypergraph over clinical records for patient-centered QA.

These methods share a limitation: they apply graph structure without domain-grounded concept annotations. Edges are built from lexical overlap, embedding similarity, or co-occurrence statistics. These are proxies; they may or may not reflect genuine biomedical relationships. None of these works systematically compares graph-based reranking with and without medical knowledge annotations. Our work provides that comparison. We compare the same hypergraph architecture with and without MeSH, entity, and PrimeKG annotations, and show that domain knowledge constraints are essential for productive graph-based reranking, while also acknowledging that a full $2{\times}2$ factorial design---with/without graph $\times$ with/without medical knowledge---would provide stronger causal isolation.

\subsection{Controlled Biomedical Knowledge}

MeSH~\citep{nlm2026mesh} is a hierarchically organized controlled vocabulary maintained by the National Library of Medicine, providing standardized annotations for all PubMed-indexed literature. PrimeKG~\citep{chandak2023primekg} integrates multimodal precision-medicine relations across diseases, proteins, drugs, and phenotypes. Both are public and widely used in biomedical information retrieval. Their use as features in controlled LTR reranking settings has not been widely studied. Our framework uses MeSH hierarchy features as the primary structured knowledge signal and PrimeKG as an auxiliary source. The design choice we want to highlight is not which knowledge base to use but how to wire structured biomedical knowledge as features inside a learning-to-rank model, so that after training we can inspect which knowledge pathways drove each ranking decision.
